\documentclass[11pt,a4paper]{letter}
\usepackage[a4paper,top=2cm,bottom=2cm,left=2.5cm,right=2.5cm]{geometry}
\usepackage[T1]{fontenc}
\usepackage[utf8]{inputenc}
\usepackage[english]{babel}
\usepackage[hidelinks]{hyperref}

\signature{Michele Ronco}
\address{Michele Ronco \\ European Commission Joint Research Centre \\ Via E. Fermi 2749, I-21027 Ispra (VA), Italy \\ \href{mailto:michele.ronco@ec.europa.eu}{michele.ronco@ec.europa.eu}}
\date{15 October 2026}

\begin{document}
\begin{letter}{The Editors \\ \emph{Genome Medicine}}

\opening{Dear Editors,}

Please find enclosed my manuscript, \emph{Chemical reachability limits mutagen-based TKI-resistance screens}, submitted as a Research Article.

The paper audits what the mutagen chemistries used in current (2024--2026) preclinical practice can and cannot produce, against two chemistry-free saturation mutagenesis atlases (Feng 2022 KRAS G12C; Oh 2026 ALK). We write a screen's expected output as $P(v) \propto S_{\mathrm{mutagen}}(v)\cdot \exp(\sigma(v))$ and find that chemistry enters as a reachability threshold rather than as a ranking weight. A cytosine-plus-adenine base-editor pair carries an exactly zero rate for four of six substitution classes, which removes eight of twelve KRAS G12C panel variants and 80--100\% of each ALK drug's chemistry-free top-10 from any such library. Two further panel variants (R68L, V8M), both reported in an accompanying \emph{in vitro} scan rather than in a patient cohort, need two nucleotide changes and are suppressed 51- to 169-fold at practical library depth. Past that threshold the class weighting does not pay its way: ENU and Mutazyme~II, whose internal class ratios differ 13-fold, give near-identical ALK top-100 recovery (61--67\%), and on no anchor does $\log S + \sigma$ beat $\sigma$ alone. We report that null as a result, and we report what the anchors can resolve ($n = 4$--11, minimum detectable gain 9--50 pp) rather than reading a non-significant $p$ as evidence of no effect. The deliverable is an enumerable, codon- and spectrum-level inventory of blind spots, auditable before a screen runs; the Discussion names the specific variants the forthcoming neladalkib and next-generation KRAS G12C screens are likely to miss.

I believe the manuscript is a fit for \emph{Genome Medicine} as a translational-methods paper: it uses genomic data to improve the preclinical screens that feed into clinical resistance panels and treatment-sequencing decisions, and its predictions are actionable both for screen designers and for clinicians interpreting preclinical-screen outputs when a patient relapses.

The work is single-authored, original, not under consideration elsewhere, and declared without competing interests. A bioR$\chi$iv preprint will accompany submission.

Suggested reviewers: James S.\ Fraser (UCSF; senior author on Estevam 2024 MET DMS); Jesse D.\ Bloom (Fred Hutchinson Cancer Center; DMS methodology); Serena Nik-Zainal (University of Cambridge; mutational signatures); Jessica J.\ Lin (Massachusetts General Hospital; clinical ALK resistance, first author on Lin 2024).

\closing{Sincerely,}

\end{letter}
\end{document}
